\chapter{Zależności zewnętrzne i poziom samowystarczalności}
\label{app:wyniki-zewnetrzne}

W skrypcie dowód oznacza rzeczywiste
wyprowadzenie ze wskazanych wcześniejszych
faktów. Jeśli korzystamy z długiego
twierdzenia spoza tekstu, podajemy
hipotezy, miejsce użycia i źródło.
Poniższy katalog ułatwia rozpoznanie
tych punktów zależności.

\section{Wyniki wejściowe i głębokie lematy}
\label{sec:katalog-wynikow}

\begin{description}[leftmargin=2.9cm,style=nextline]
\item[Niezmienniczość obszaru.]
  Ciągła injekcja z otwartego
  $U\subset\R^n$ do $\R^n$
  ma otwarty obraz. Użyta w
  Rozdziale~\ref{ch:rozmaitosci}
  do niezależności wymiaru i
  brzegu od map.
  Źródło:
  \href{https://pi.math.cornell.edu/~hatcher/AT/AT.pdf}
  {A.\ Hatcher, \emph{Algebraic Topology},
  twierdzenie 2B.3}.
  Alternatywny dowód samej
  niezmienniczości wymiaru rozmaitości
  przez homologię lokalną jest poniżej.

\item[Twierdzenie o funkcji odwrotnej.]
  Dla $C^1$ odwzorowania między
  otwartymi podzbiorami $\R^n$
  z odwracalną pochodną w punkcie
  istnieje lokalna odwrotność $C^1$;
  dla danych $C^\infty$ jest ona
  $C^\infty$. Użyte przy mapach,
  zanurzeniach i funkcjach Morse'a.
  Źródło z dowodem:
  \href{https://www.ocw.mit.edu/courses/18-101-analysis-ii-fall-2005/d2c8fc5fe57cd59efc99caa048983d4c_lecture6.pdf}
  {V.\ Guillemin, \emph{Analysis II},
  wykład 6}.

\item[Miara Hausdorffa w przestrzeni euklidesowej.]
  Identyfikacja znormalizowanej miary Hausdorffa
  $\mathcal H^k$ na $\R^k$ z miarą Lebesgue'a,
  użyta w Rozdziale~\ref{ch:metryka-riemannowska},
  korzysta z nierówności izodiametrycznej i z twierdzenia
  Vitalego o pokryciu. Sam rachunek normalizacji oraz
  zastosowanie w mapach rozmaitości są w rozdziale,
  natomiast te dwa fakty miarowe pozostają zewnętrzne.
  Źródło:
  \href{https://www.math.cmu.edu/~gautam/sj/teaching/2022-23/720-measure/pdfs/measure.pdf}
  {G.\ Iyer, \emph{Lecture Notes on Measure Theory},
  dodatek A}.

\item[Triangulacja gładka.]
  Każda gładka rozmaitość
  drugoprzeliczalna dopuszcza
  triangulację zgodną ze
  strukturą gładką; w przypadku
  zwartym może być skończona.
  Używamy jej do przejścia od
  geometrii do skończonych
  kompleksów łańcuchowych.
  Źródło:
  \href{https://www.ams.org/bull/1935-41-08/S0002-9904-1935-06140-3/S0002-9904-1935-06140-3.pdf}
  {S.~S.\ Cairns, \emph{Triangulation of the
  Manifold of Class One}, §1--2}, oraz
  \href{https://doi.org/10.2307/1968861}
  {J.~H.~C.\ Whitehead, \emph{On $C^1$-complexes}}.
  Pełny dowód nie jest krótkim
  uzupełnieniem tego skryptu.

\item[Kohomologia Čech i singularna.]
  Każde pokrycie rozmaitości ma
  rozdrobnienie będące dobrym pokryciem.
  Dobre pokrycia tworzą więc rodzinę
  kofinalną w definicji kohomologii
  Čecha. Dla dobrego pokrycia kompleks
  Čecha oblicza kohomologię nerwu,
  a lemat o nerwie utożsamia jego
  typ homotopijny z rozmaitością.
  Dla współczynników w dyskretnej
  grupie abelowej otrzymujemy stąd
  izomorfizm kohomologii Čecha
  i singularnej. Interpretacja
  kocyklu przejść struktury
  spinowej w Rozdziale~\ref{ch:grass-spin}
  używa tego porównania;
  tam, gdzie wystarcza stopień
  pierwszy, podano też argument
  przez monodromię pętli.
  Istnienie dobrych pokryć opisują
  \href{https://doi.org/10.1007/978-1-4757-3951-0}
  {R.\ Bott i L.~W.\ Tu, \emph{Differential Forms in
  Algebraic Topology}, §5}. Wybierając użyte tam dostatecznie
  małe geodezyjnie wypukłe kule podporządkowane zadanemu pokryciu,
  otrzymujemy dobre rozdrobnienie; to daje potrzebną kofinalność.
  Pozostałe składniki porównania:
  \href{https://pi.math.cornell.edu/~hatcher/AT/AT.pdf}
  {A.\ Hatcher, \emph{Algebraic Topology},
  wniosek 4G.3 o nerwie dobrego pokrycia
  oraz §2.1 i §3.1 o kohomologii
  symplicjalnej i singularnej}.

\item[Trik Whitneya i ruchy uchwytów.]
  Niech $P^r,Q^s$ będą domkniętymi, poprzecznie
  przecinającymi się gładkimi podrozmaitościami
  $V^{r+s}$ bez brzegu; orientujemy $P$ i wiązkę
  normalną $Q$. Usunięcie punktów przecięcia $u,v$
  o przeciwnych znakach wymaga $r\geq2$, $s\geq3$
  i $r+s\geq5$. Łuki łączące $u,v$ w $P$ i $Q$
  mają wnętrzami omijać pozostałe przecięcia, a utworzony
  przez nie okrąg ma być ściągalny w $V$. Przy $r=2$
  zakładamy ponadto injektywność
  $\pi_1(V\setminus Q)\to\pi_1(V)$. Konstrukcja daje
  odpowiednio obramowany dysk Whitneya i izotopię
  usuwającą tę parę przecięć.
  Użyte w
  Rozdziałach~\ref{ch:hcobordyzm}--\ref{ch:przeszkoda-chirurgiczna}.
  Źródło warunków i konstrukcji:
  \href{https://www.math.columbia.edu/~jmorgan/Lecture_IIIA_hcobordism_Contd.pdf}
  {J.\ Morgan, \emph{The $h$-Cobordism Theorem},
  wykład IIIA}, oraz
  \href{https://webhomes.maths.ed.ac.uk/~v1ranick/surgery/hcobord.pdf}
  {J.\ Milnor, \emph{Lectures on the $h$-Cobordism Theorem},
  twierdzenie 6.6 i lematy 6.7, 6.13}.

\item[Dwie chirurgii okręgu w $S^4$.]
  Dla standardowego okręgu dwa zorientowane obramowania
  dają produktową i nietrywialną wiązkę $S^2$ nad $S^2$.
  Rozdział~\ref{ch:chirurgia-geometryczna-21} oblicza ich formy
  przecięcia; identyfikacja nietrywialnej wiązki z
  $\mathbb{CP}^2\#\overline{\mathbb{CP}^2}$ pozostaje wejściem z
  \href{https://arxiv.org/abs/1502.06834}
  {K.\ Larsona, \emph{Surgery on tori in the 4-sphere},
  dowód propozycji 4.3} oraz
  \href{https://www.mathematik.hu-berlin.de/~wendl/pub/rationalRuled.pdf}
  {C.\ Wendla, \emph{Holomorphic Curves in Low Dimensions},
  przykład 1.11 i uwaga 1.13}.
  Nie jest to zastosowanie wysokowymiarowego
  twierdzenia Walla.

\item[Torsja Whiteheada i $s$-kobordyzm.]
  Dla równoważności homotopijnej spójnych skończonych
  kompleksów CW zerowanie torsji Whiteheada jest równoważne
  prostocie tej równoważności. Wynik ten stosujemy w
  Rozdziale~\ref{ch:s-kobordyzm}; źródło:
  \href{https://webhomes.maths.ed.ac.uk/~v1ranick/papers/owsem.pdf}
  {M.\ Kreck i W.\ Lück, \emph{The Novikov Conjecture},
  twierdzenie 6.15(1)}. Drugim wejściem jest utożsamienie
  torsji uporządkowanego kompleksu Morse'a z torsją
  komórkową zgodnej gładkiej triangulacji oraz niezależność
  tej klasy od triangulacji; zob.
  \href{https://webhomes.maths.ed.ac.uk/~v1ranick/papers/milnorwh.pdf}
  {J.\ Milnor, \emph{Whitehead Torsion}, lemat 9.1
  i twierdzenie 9.3}. Po tych wejściach rozdział dowodzi,
  że dla spójnego zwartego gładkiego $h$-kobordyzmu
  $(W^m;N_-,N_+)$ o niepustych, spójnych i zamkniętych
  brzegach, $m\geq6$, zachodzi
  $\tau(W,N_-)=0$ wtedy i tylko wtedy, gdy
  $W\cong N_-\times I$ relatywnie do $N_-$;
  orientowalność nie jest założeniem.
  Przykład 6.6 Milnora oblicza osobno
  $\operatorname{Wh}(C_5)\cong\Z$ z generatorem
  $[1-t-t^{-1}]$; używamy go w
  Przykładzie~\ref{prz:Tate-C5-22}.

\item[Twierdzenia Walla.]
  Niech $(f,b):M^n\to X^n$ będzie normalnym
  odwzorowaniem stopnia jeden między spójnymi,
  zamkniętymi, gładkimi rozmaitościami, $n\geq5$,
  z charakterem orientacji $w$ na
  $\pi=\pi_1(X)$. Przeszkoda
  $\theta(f,b)\in L_n^s(\Z[\pi],w)$ znika wtedy
  i tylko wtedy, gdy $(f,b)$ jest normalnie
  kobordantne z prostą równoważnością homotopijną.
  W wersji zwartej z brzegiem odwzorowanie brzegowe
  ma być prostą równoważnością, a kobordyzm i chirurgie
  wykonuje się relatywnie do brzegu. Przypadek parzysty
  $n=2k\geq6$ daje twierdzenie 5.6 Walla, a przypadek
  nieparzysty $n=2k+1\geq5$, obejmujący wymiar pięć,
  daje twierdzenie 6.4; nie należy go mylić z nieparzystym
  twierdzeniem realizacyjnym 6.5, którego zakres zaczyna
  się w wymiarze siedem.

  Rozdział~\ref{ch:przeszkoda-chirurgiczna}
  rozpisuje rachunek przecięć i konstrukcję
  bazy hiperbolicznej nad $\Z[\pi]$,
  chirurgie poniżej środka, ciągi dokładne jąder,
  dualność oraz część ruchów formacji. Osobnymi
  wejściami pozostają przygotowanie kobordyzmu,
  bazowość i prosta dualność jego jąder,
  realizacja względnych immersji, geometryczne lematy
  Walla oraz pełna realizacja ruchów nieparzystych.

  Wersja z jedynie zwykłą równoważnością na brzegu ma
  przeszkodę $\theta^h(f,b)\in L_n^h(\Z[\pi],w)$.
  Homomorfizm zapominający prostotę
  $L_n^s\to L_n^h$ i jego związek z grupą
  Whiteheada opisuje ciąg dokładny Rothenberga z grupami
  Tate'a $\widehat H^*(\Z/2;\operatorname{Wh}(\pi))$.
  Te twierdzenia są zewnętrzne; zob.
  \href{https://www.maths.ed.ac.uk/~v1ranick/surgery/ranicki.pdf}
  {A.\ Ranicki, \emph{An Introduction to Algebraic Surgery},
  wstęp i przykład 2.16},
  \href{https://webhomes.maths.ed.ac.uk/~v1ranick/books/scm.pdf}
  {Walla, §17D}.

  Osobnym wejściem jest rozszczepienie Shanesona--Ranickiego
  $L_n^s(\Z[G\times\Z],w\oplus1)
   \cong L_n^s(\Z[G],w)\oplus L_{n-1}^h(\Z[G],w)$
  dla dodatkowego generatora zachowującego orientację.
  W przypadku $G=1$ prowadzi ono do tabeli
  $L_n^s(\Z[t,t^{-1}],1)$ w
  Sekcji~\ref{sec:L-pi1-Z-22}. Źródła:
  \href{https://webhomes.maths.ed.ac.uk/~v1ranick/papers/shan.pdf}
  {J.\ Shaneson, \emph{Wall's Surgery Obstruction Groups
  for $G\times\Z$}, twierdzenie 5.1, wariant geometryczny
  dla $G=\pi_1(K^{n-2})$, $w=w_1(K)$, $n\geq7$
  i zamkniętej gładkiej lub PL rozmaitości $K$} oraz
  \href{https://webhomes.maths.ed.ac.uk/~v1ranick/papers/splitkl.pdf}
  {A.\ Ranicki, \emph{Algebraic and Geometric Splittings
  of the K- and L-Groups of Polynomial Extensions}, §6,
  wariant algebraiczny dla dowolnej grupy}.

  Dla spójnej zamkniętej gładkiej $X^n$, $n\geq5$,
  identyfikacja $\mathcal N_{\mathrm{DIFF}}(X)\cong[X,G/O]$,
  działanie $L_{n+1}^s$ na prostym zbiorze struktur i ciąg
  chirurgii
  \[
   [\Sigma X_+,G/O]\longrightarrow
   L_{n+1}^s(\Z[\pi],w)
   \longrightarrow\mathcal S^s_{\mathrm{DIFF}}(X)
   \longrightarrow[X,G/O]\longrightarrow
   L_n^s(\Z[\pi],w)
  \]
  są kolejnymi zewnętrznymi wynikami Walla użytymi w
  Rozdziale~\ref{ch:ciag-chirurgii}. Realizacja elementów
  $L_{n+1}^s$ korzysta z twierdzenia 5.8 dla parzystego
  $n+1\geq6$ oraz z twierdzenia 6.5 dla nieparzystego
  $n+1\geq7$; przy $n=5$ używa się pierwszego z nich.
  Główne źródło:
  \href{https://webhomes.maths.ed.ac.uk/~v1ranick/books/scm.pdf}
  {C.~T.~C. Wall, \emph{Surgery on Compact
  Manifolds}, §§1--2, 5--6, 10 i 17D,
  zwłaszcza twierdzenia 1.2, 5.2, 5.6, 5.8,
  6.4--6.5, lemat 5.7 oraz wyniki 10.1--10.8}.
  Rozdziały~\ref{ch:przeszkoda-chirurgiczna}
  oraz~\ref{ch:ciag-chirurgii} rozpisują zastosowania,
  lecz nie zastępują tych wyników geometrycznych.

\item[Triangulacja powierzchni.]
  Dowód klasyfikacji zwartych powierzchni
  w Twierdzeniu~\ref{thm:classification-surfaces-24}
  korzysta z istnienia skończonej triangulacji
  topologicznej powierzchni; źródło:
  \href{https://pi.math.cornell.edu/~hatcher/Papers/TorusTrick.pdf}
  {A.\ Hatcher, \emph{The Kirby Torus Trick for Surfaces}}.

\item[Klasyfikacja Bardena w wymiarze pięć.]
  Jeśli $P,Q$ są zamkniętymi, zorientowanymi, gładkimi
  i jednospójnymi $5$-rozmaitościami, to izomorfizm
  $H_2(P;\Z)\to H_2(Q;\Z)$ zachowujący formę wiązania
  na torsji i klasę $w_2$ jest realizowany przez
  dyfeomorfizm zachowujący orientację. Jest to
  \href{https://www.sas.rochester.edu/mth/sites/doug-ravenel/otherpapers/barden.pdf}
  {D.\ Barden, \emph{Simply Connected
  Five-Manifolds}, twierdzenie 2.2}, zapisane jako
  Twierdzenie~\ref{thm:Barden-5-24}. Kanoniczna
  klasyfikacja modeli w
  Twierdzeniu~\ref{thm:barden-kanoniczna-24}
  jest twierdzeniem 2.3 tej samej pracy.
  Rozdział~\ref{ch:poincare-2-high} rekonstruuje argument,
  ale przyjmuje z zewnątrz obliczenia grup kobordyzmu
  $\Omega^{SO}_5$ i $\Omega^{Spin}_5$ Thoma--Nowikowa,
  uogólnione twierdzenie Hurewicza Whiteheada,
  reprezentowanie klas przez zanurzenia Haefligera oraz
  klasyfikację form wiązania Walla--Bardena. Dokładne
  miejsca w źródłach to Barden, §§0.3--0.4,
  twierdzenia 2.2--2.3 oraz lematy 3.1 i 4.1--4.4, a także
  \href{https://webhomes.maths.ed.ac.uk/~v1ranick/papers/kill.pdf}
  {C.~T.~C. Wall, \emph{Killing the Middle Homotopy Groups
  of Odd Dimensional Manifolds}, stwierdzenia 1--2
  i lemat 4(ii)}.

\item[Sygnatura i sfery Milnora.]
  Dla zamkniętej zorientowanej gładkiej
  $8$-rozmaitości $C$ używamy
  $\operatorname{sign}(C)=
  \langle(7p_2-p_1^2)/45,[C]\rangle$.
  Źródło:
  \href{https://webhomes.maths.ed.ac.uk/~v1ranick/papers/hirzrem.pdf}
  {F.\ Hirzebruch, \emph{The Signature Theorem},
  wzór na $L_2$}.
  Wzory charakterystyczne wiązek nad $S^4$ pochodzą z
  \href{https://sites.math.rutgers.edu/~feehan/teaching/math866/milnor7sphere.pdf}
  {J.\ Milnor, \emph{On Manifolds Homeomorphic
  to the 7-Sphere}, lemat 3}.
  Niezależność niezmiennika Eellsa--Kuipera
  $\mu\in\Q/\Z$ od spinowego wypełnienia korzysta ponadto
  z całkowitości indeksu $\widehat A$ operatora Diraca
  na zamkniętej spinowej $8$-rozmaitości. Źródłem wzoru
  i jego zastosowania jest
  \href{https://webhomes.maths.ed.ac.uk/~v1ranick/papers/eelskuip.pdf}
  {J.\ Eells i N.\ Kuiper, \emph{An Invariant for Certain
  Smooth Manifolds}, §6}. Klasyfikacja w
  Rozdziale~\ref{ch:sfery-Milnora} używa jako osobnych
  wyników zewnętrznych równości
  $\Theta_7=bP_8\cong\Z/28$ z
  \href{https://www.sas.rochester.edu/mth/sites/doug-ravenel/otherpapers/kervaire-milnor.pdf}
  {M.\ Kervaire'a i J.\ Milnora, \emph{Groups of Homotopy
  Spheres I}, §4 i §§7--8}, oraz zupełności $\mu$ dla
  zorientowanych gładkich sfer homeomorficznych z $S^7$
  z §6 pracy Eellsa--Kuipera.

\item[Holonomia i globalny rozkład de Rhama.]
  Dla koneksji Leviego-Civity twierdzenie
  Ambrose'a--Singera mówi, że algebra Liego holonomii
  ograniczonej jest rozpięta przez operatory
  $P_\gamma^{-1}R_q(u,v)P_\gamma$; obejmuje także fakt,
  że holonomia ograniczona jest grupą Liego. Wynik ten,
  użyty w Rozdziale~\ref{ch:krzywizna}, pochodzi z
  \href{https://doi.org/10.1090/S0002-9947-1953-0063739-1}
  {W.\ Ambrose'a i I.~M.\ Singera,
  \emph{A theorem on holonomy} (1953)}.
  W tym samym rozdziale udowodniono lokalną postać
  iloczynu, lecz globalny krok jest osobnym twierdzeniem:
  spójna, zupełna i jednospójna rozmaitość riemannowska
  bez brzegu, której pełna holonomia jest redukowalna,
  rozkłada się izometrycznie na czynnik euklidesowy
  i czynniki nieredukowalne, jednoznacznie z dokładnością
  do izometrii i przestawienia czynników. Źródła:
  \href{https://doi.org/10.1007/BF02564308}
  {G.\ de Rham, \emph{Sur la réductibilité d'un espace
  de Riemann} (1952)} oraz
  \href{https://doi.org/10.1215/ijm/1256059674}
  {H.\ Wu, \emph{On the de Rham decomposition theorem}
  (1964)}.

\item[Normalizacja topologiczna Cherna--Weila.]
  Rozdział~\ref{ch:chern-weil} dowodzi zamkniętości
  form Cherna--Weila i wzoru transgresji. Zewnętrznym
  wejściem pozostaje zgodność formy wyznaczonej przez
  $\det(I+i\Omega/(2\pi))$ z rzeczywistym obrazem
  topologicznej klasy Cherna, w szczególności $c_2$,
  oraz tożsamość $p_1(E)=-c_2(E\otimes_\R\C)$.
  Źródła:
  \href{https://pi.math.cornell.edu/~hatcher/VBKT/VB.pdf}
  {A.\ Hatcher, \emph{Vector Bundles and K-Theory},
  §3.1--3.2}, oraz
  \href{https://webusers.imj-prg.fr/~olivier.biquard/m2dg2022.pdf}
  {O.\ Biquard, \emph{Differential and Riemannian Geometry},
  §16.c, wzory (III.39) i (III.41)}.

\item[Analityczna teoria Hodge'a.]
  Dla laplasjanu Hodge'a na formach zwartej
  rozmaitości riemannowskiej bez brzegu
  jądro jest skończenie wymiarowe, a każda
  gładka forma prostopadła do jądra ma
  jednoznaczny gładki przeciwobraz
  prostopadły do jądra. Jest to
  Twierdzenie~\ref{thm:hodge-analytic-input},
  użyte w Rozdziale~\ref{ch:hodge-bochner}
  do istnienia rozkładu; jego konsekwencje
  algebraiczne i wzór Bochnera tam dowodzimy.
  Źródło:
  \href{https://doi.org/10.2307/1969552}
  {K.\ Kodaira, \emph{Harmonic Fields in
  Riemannian Manifolds} (1949)}.

\item[Analityczna teoria potoku Ricciego.]
  Gładki ściśle paraboliczny układ
  quasiliniowy drugiego rzędu na zwartej
  rozmaitości bez brzegu ma dla gładkich
  danych lokalnie w czasie jednoznaczne
  gładkie rozwiązanie; używamy także
  odpowiedniej wersji dla równania ciepła
  odwzorowań z metryką dziedziny zależną
  od czasu. W
  Twierdzeniu~\ref{thm:ricci-short-time}
  te fakty stosujemy po poprawce DeTurcka.
  Wyprowadzenie głównego symbolu i cofnięcie
  metryki przez dyfeomorfizmy są w rozdziale;
  same oszacowania paraboliczne są zewnętrzne.
  Źródło twierdzenia o krótkim czasie:
  \href{https://projecteuclid.org/journals/journal-of-differential-geometry/volume-17/issue-2/Three-manifolds-with-positive-Ricci-curvature/10.4310/jdg/1214436922.full}
  {R.\ S.\ Hamilton, \emph{Three-manifolds with
  positive Ricci curvature} (1982), \S\,4}.
  Źródło konstrukcji i równania cechowania:
  \href{https://projecteuclid.org/journals/journal-of-differential-geometry/volume-18/issue-1/Deforming-metrics-in-the-direction-of-their-Ricci-tensors/10.4310/jdg/1214509286.full}
  {D.\ DeTurck, \emph{Deforming Metrics in the
  Direction of Their Ricci Tensors} (1983)}.

\item[Wygładzanie rozmaitości trójwymiarowych.]
  Każda topologiczna rozmaitość trójwymiarowa dopuszcza
  zgodną strukturę gładką. Ten wynik Moise'a
  pozwala przejść od gładkiej postaci
  Twierdzenia~\ref{thm:poincare-trzy}
  do topologicznej postaci we
  Wniosku~\ref{cor:poincare-trzy-topologiczna}.
  Źródło: \href{https://msp.org/gtm/2015/19-1/gtm-v19-n1-p08-s.pdf}
  {M. Aschenbrenner, S. Friedl i H. Wilton,
  \emph{Decision Problems for 3-Manifolds and Their
  Fundamental Groups}, twierdzenie 3.1};
  \href{https://doi.org/10.2307/1969769}
  {E.~E. Moise, \emph{Affine Structures in
  3-Manifolds V} (1952)}.

\item[Miara Haara i ilorazy grup Liego.]
  Na każdej zwartej grupie Liego istnieje znormalizowana
  miara probabilistyczna niezmiennicza przy mnożeniu
  z lewej i z prawej strony. W
  Rozdziale~\ref{ch:geometria-kleina} służy ona do
  uśrednienia metryki względem zwartego stabilizatora.
  Źródło:
  \href{https://doi.org/10.2307/1968346}
  {A.\ Haar, \emph{Der Massbegriff in der Theorie der
  kontinuierlichen Gruppen} (1933)}.
  Dla domkniętej podgrupy $H$ skończenie wymiarowej
  grupy Liego $G$ istnienie gładkiej struktury na
  $G/H$ i submersyjność rzutowania są wejściem
  zewnętrznym w
  Twierdzeniu~\ref{thm:klein-iloraz-liego}.
  Dowód identyfikacji orbit z ilorazami oraz
  konkretne przykłady są w
  Rozdziale~\ref{ch:geometria-kleina}.
  Źródło: J.\ M.\ Lee, \emph{Introduction to Smooth
  Manifolds}, rozdziały~20--21.

\item[Potok Ricciego z chirurgią i wygaśnięcie.]
  Dla dowolnej gładkiej metryki na zamkniętej
  zorientowanej $3$-rozmaitości kontrolę osobliwości
  i kontynuację potoku przez chirurgię dostarczają
  wyniki Perelmana. Używana wersja twierdzenia o
  wygaśnięciu zakłada, że $\pi_1(M)$ jest wolnym
  iloczynem skończenie wielu grup skończonych i kopii
  $\Z$; równoważnie, w rozkładzie pierwszym nie ma
  składników asferycznych. Wówczas potok z chirurgią
  wygasa w skończonym czasie.
  W Rozdziale~\ref{ch:poincare-trzy} te twierdzenia
  oraz topologiczny opis odrzucanych składników
  są jawnymi wejściami; własne rachunki i wnioski
  nie zastępują ich dowodów. Źródła:
  \href{https://arxiv.org/pdf/math/0211159}
  {G.\ Perelman, \emph{The entropy formula for the Ricci flow} (2002)},
  \href{https://arxiv.org/pdf/math/0303109}
  {\emph{Ricci flow with surgery on three-manifolds} (2003)},
  \href{https://arxiv.org/pdf/math/0307245}
  {\emph{Finite extinction time} (2003)};
  szczegółowe opracowanie:
  \href{https://www.claymath.org/wp-content/uploads/2022/03/Ricci-pdf.pdf}
  {J.\ Morgan i G.\ Tian, \emph{Ricci Flow and the Poincaré Conjecture}}.

\item[Geometryzacja rozmaitości trójwymiarowych.]
  Rozkład pierwszy, cięcie wzdłuż istotnych torusów
  i twierdzenie geometryzacyjne dla zamkniętych
  orientowalnych $3$-rozmaitości są wynikami zewnętrznymi
  w Rozdziale~\ref{ch:geometrizacja-thurstona}.
  Opis ośmiu maksymalnych geometrii modelowych także
  korzysta z klasyfikacji zewnętrznej; przykłady i rachunki
  w rozdziale ilustrują modele, lecz jej nie dowodzą.
  Źródła:
  \href{https://library.slmath.org/nonmsri/gt3m/PDF/Thurston-gt3m.pdf}
  {W.\ Thurston, \emph{The Geometry and Topology of Three-Manifolds}},
  \href{https://sschleimer.warwick.ac.uk/Teach/2021_3MFDS/References/1983the_geometries_of_three-manifolds.pdf}
  {P.\ Scott, \emph{The Geometries of 3-Manifolds} (1983)}
  oraz \href{https://www.claymath.org/wp-content/uploads/2022/03/Geom.pdf}
  {J.\ Morgan i G.\ Tian, \emph{The Geometrization Conjecture}}.

\item[Topologiczna dualność i wygładzanie PL w wymiarze cztery.]
  W Rozdziale~\ref{ch:freedman-formy} formę przecięcia
  zamkniętej zorientowanej rozmaitości TOP opisujemy przez topologiczną
  dualność Poincarégo. Używamy także wyniku, że
  struktura PL w wymiarze cztery dopuszcza zgodne
  wygładzenie. Rachunki macierzy uchwytów i formy $E_8$
  są w rozdziale; te dwa fakty kategorialne pozostają
  wejściami zewnętrznymi. Źródło:
  \href{https://archive.mpim-bonn.mpg.de/4789/2/FreedmanQuinn-TopologyOf4Manifolds-Reformatted2013.pdf}
  {M.\ H.\ Freedman i F.\ Quinn, \emph{Topology of
  4-Manifolds}, \S\,8.7 i \S\,10.2}.

\item[Zanurzanie dysków i uchwyty Cassona.]
  W Rozdziale~\ref{ch:uchwyty-cassona} dowodzimy
  własności skończonych pięter, ale przejście do
  parami rozłącznych lokalnie płaskich dysków jest
  wejściem zewnętrznym. Użyta wersja dotyczy zwartej,
  spójnej, zorientowanej gładkiej $4$-rozmaitości
  z brzegiem i właściwych immersji obramowanych dysków
  o parami rozłącznie zanurzonych brzegach z kołnierzami.
  Zakłada obramowane immersyjne sfery $B_j$ z
  $\lambda(A_i,B_j)=\delta_{ij}$,
  $\lambda(B_i,B_j)=0$ i $\mu(B_j)=0$ nad
  $\Z[\pi_1(M)]$ oraz grupę podstawową wielostopniowo
  skończoną lub cykliczną. Konkluzja daje dyski
  topologicznie zanurzone z zachowanymi obramowanymi
  brzegami i geometrycznie dualnymi sferami:
  Twierdzenie~\ref{thm:casson-disk-embedding},
  \href{https://archive.mpim-bonn.mpg.de/4789/2/FreedmanQuinn-TopologyOf4Manifolds-Reformatted2013.pdf}
  {Freedman--Quinn, \emph{Topology of 4-Manifolds},
  Twierdzenie 5.2}. Konstrukcję geometrycznie
  dualnych sfer w konkluzji uzupełnia
  \href{https://arxiv.org/pdf/2006.05209}
  {M.\ Powell, A.\ Ray i P.\ Teichner,
  \emph{The 4-dimensional disc embedding theorem
  and dual spheres}, Twierdzenie A}; w pierwotnym
  dowodzie Freedmana--Quinna brakowało tego kroku.
  Standardowość topologiczna
  uchwytu Cassona jako pary to osobne wejście:
  Twierdzenie~\ref{thm:casson-homeomorphism},
  \href{https://www.maths.gla.ac.uk/~mpowell/1982_The%20topology%20of%20four-dimensional%20manifolds.pdf}
  {Freedman, \emph{The topology of four-dimensional
  manifolds} (1982), Twierdzenie 1.1}.

\item[Klasyfikacja Freedmana i niezmiennik Kirby'ego--Siebenmanna.]
  Twierdzenie~\ref{thm:freedman-classification}
  o zamkniętych zorientowanych jednospójnych
  rozmaitościach TOP, warunek
  $\operatorname{ks}\equiv\sigma/8\pmod2$ dla
  form parzystych oraz
  dokładne kryterium stabilnego wygładzenia
  $\operatorname{ks}(X)=0\Longleftrightarrow
  X\times\R$ jest wygładzalna
  są wejściami zewnętrznymi w
  Rozdziale~\ref{ch:freedman-klasyfikacja}.
  Kryterium to nie orzeka samo o wygładzalności $X$.
  Źródła:
  \href{https://archive.mpim-bonn.mpg.de/4789/2/FreedmanQuinn-TopologyOf4Manifolds-Reformatted2013.pdf}
  {Freedman--Quinn, \emph{Topology of 4-Manifolds},
  \S\,8.7--8.8 i \S\,10.1--10.2} oraz
  \href{https://www.maths.gla.ac.uk/~mpowell/1982_The%20topology%20of%20four-dimensional%20manifolds.pdf}
  {Freedman, \emph{The topology of four-dimensional
  manifolds} (1982), Twierdzenie 1.5}.

\item[Twierdzenie Donaldsona i analiza równań ASD.]
  Diagonalizacja określonej całkowitej formy przecięcia
  gładkiej, zamkniętej, jednospójnej $4$-rozmaitości
  jest zewnętrznym Twierdzeniem~\ref{thm:donaldson-diagonalizacja}.
  Rachunek energii w Rozdziale~\ref{ch:donaldson} jest własny,
  lecz zwartość przestrzeni moduli, analiza rozwiązań
  redukowalnych i koncentracji krzywizny nie są tu dowiedzione.
  Źródło:
  \href{https://doi.org/10.4310/jdg/1214437665}
  {S.~K.~Donaldson, \emph{An application of gauge theory
  to four-dimensional topology}, J.~Differential Geom. 18
  (1983), twierdzenie A}.

\item[Twierdzenie Rokhlina.]
  Podzielność sygnatury zamkniętej gładkiej
  spinowej $4$-rozmaitości przez $16$, użyta do
  porównania form $E_8$ i $E_8\oplus E_8$ w
  Rozdziale~\ref{ch:donaldson}, jest wynikiem
  zewnętrznym. Zob.
  \href{https://arxiv.org/pdf/2011.12418}
  {M.~R.~Klug, \emph{A relative version of Rochlin's
  theorem}, wprowadzenie}.

\item[Egzotyczne struktury gładkie na $\R^4$.]
  Istnienie początkowej zwartej przeszkody, realizacja
  śladu topologicznie plastrowalnego splotu oraz
  lemat wytwarzający następną przeszkodę są zewnętrznymi
  Twierdzeniami~\ref{thm:egzotyczne-r4-X},
  \ref{thm:egzotyczne-r4-slad} i
  \ref{thm:egzotyczne-r4-lemat}. Indukcja i rozróżnienie
  typów gładkich są rozpisane w
  Rozdziale~\ref{ch:egzotyczne-r4}.
  Źródło:
  \href{https://doi.org/10.4310/jdg/1214439566}
  {R.~E.~Gompf, \emph{An infinite set of exotic
  $\R^4$'s}, J.~Differential Geom. 21 (1985),
  lematy 1.1--1.2 i twierdzenie 1.3}.
  Dowód realizacji śladu splotu korzysta także z
  zewnętrznych twierdzeń Quinna o wygładzaniu
  niezwartych dopełnień w wymiarze cztery i o stabilnych
  homeomorfizmach; zob.
  \href{https://www.maths.gla.ac.uk/~mpowell/1982_Ends%20of%20maps%20III.pdf}
  {F.\ Quinn, \emph{Ends of maps, III: dimensions 4 and 5}
  (1982)}.
  Istnienie kontinuum małych egzotycznych $\R^4$
  w Twierdzeniu~\ref{thm:egzotyczne-r4-male}
  jest osobnym wynikiem zewnętrznym:
  \href{https://doi.org/10.4310/jdg/1214447810}
  {S.~De Michelis i M.~H.~Freedman,
  \emph{Uncountably many exotic $\R^4$'s in
  standard 4-space}, J.~Differential Geom. 35
  (1992), 219--254}.
\end{description}
